\subsubsection{Матрична форма двовимірної барицентричної інтерполяції}

Хоча класичний алгебраїчний запис двовимірної барицентричної формули~\eqref{eq:barycentric_2d_general} сам по собі забезпечує високу чисельну стійкість та алгоритмічну ефективність, наступним логічним етапом оптимізації розробленого методу є перехід до його матричного представлення.

Такий перехід зумовлений специфікою обробки цифрових зображень. Оскільки вагові коефіцієнти в барицентричній інтерполяції залежать виключно від просторових координат (вузлів сітки та цільових точок) і не залежать від значень інтенсивності пікселів, розрахунок геометричних ваг та власне інтерполяцію кольору доцільно розділити. Для забезпечення максимальної швидкодії при обробці зображень високої роздільної здатності математичну модель необхідно адаптувати під архітектуру сучасних обчислювальних систем.

У цьому підрозділі виконується математичний перехід від базового запису до матричної форми. Цей етап оптимізації дозволяє звести задачу двовимірного масштабування до двох послідовних матричних множень. З інженерної точки зору це є критично важливим кроком, оскільки матрична форма забезпечує кеш-дружній (послідовний) доступ до пам'яті, уможливлює розпаралелювання обчислень та відкриває шлях до глибокої апаратної оптимізації алгоритму за допомогою векторизованих інструкцій процесора.

\paragraph{Крок~1. Факторизація ядра.}
Розглянемо базову формулу для довільної цільової точки $(x, y)$:
\begin{equation}\label{eq:m_original}
	L_{n,m}(x, y) = \frac{
		\sideset{}{'}\sum\limits_{i=0}^{n}
		\sideset{}{'}\sum\limits_{j=0}^{m}
		\dfrac{(-1)^{i+j}}{(x - x_i)(y - y_j)}\, f_{i,j}
	}{
		\sideset{}{'}\sum\limits_{i=0}^{n}
		\sideset{}{'}\sum\limits_{j=0}^{m}
		\dfrac{(-1)^{i+j}}{(x - x_i)(y - y_j)}
	}.
\end{equation}

Введемо ненормалізовані вагові коефіцієнти, розділивши ядро на множники, що залежать лише від однієї змінної. Скористаємося тотожностями $(-1)^{i+j} = (-1)^i \cdot (-1)^j$ та $\frac{1}{(x-x_i)(y-y_j)} = \frac{1}{x-x_i} \cdot \frac{1}{y-y_j}$:
\begin{equation}\label{eq:m_uv}
	u_i(x) \stackrel{\text{def}}{=} \frac{w_i}{x - x_i}, \qquad
	v_j(y) \stackrel{\text{def}}{=} \frac{w_j}{y - y_j},
\end{equation}
де $w_k = (-1)^k$ з половинними значеннями на краях: $w_0 = \tfrac{1}{2}$, $w_n = \tfrac{(-1)^n}{2}$ (аналогічно для осі~$y$). Позначимо чисельник і знаменник формули~\eqref{eq:m_original} як $N$ та $D$ відповідно. Тоді:
\begin{equation}\label{eq:m_ND}
	N = \sum\limits_{i=0}^{n}\sum\limits_{j=0}^{m} u_i(x)\, v_j(y)\, f_{i,j}, \qquad
	D = \sum\limits_{i=0}^{n}\sum\limits_{j=0}^{m} u_i(x)\, v_j(y).
\end{equation}

\paragraph{Крок~2. Факторизація знаменника.}
Оскільки знаменник не містить множника $f_{i,j}$, що зв'язує індекси, подвійна сума розпадається у добуток двох одинарних:
\begin{equation}\label{eq:m_D_factor}
	D = \left(\sum\limits_{i=0}^{n} u_i(x)\right) \cdot \left(\sum\limits_{j=0}^{m} v_j(y)\right)
	\stackrel{\text{def}}{=} S_u(x) \cdot S_v(y).
\end{equation}

\paragraph{Крок~3. Розподіл ділення.}
Перегрупуємо чисельник - винесемо за внутрішню суму множники, що не залежать від~$j$:
\begin{equation}\label{eq:m_N_regroup}
	N = \sum\limits_{i=0}^{n} u_i(x) \cdot \left[\sum\limits_{j=0}^{m} v_j(y)\, f_{i,j}\right].
\end{equation}
Підставимо у вираз $L = N/D$:
\begin{equation}\label{eq:m_L_step}
	L = \frac{
		\sum\limits_{i=0}^{n} u_i(x) \cdot
		\left[\sum\limits_{j=0}^{m} v_j(y)\, f_{i,j}\right]
	}{S_u(x) \cdot S_v(y)}.
\end{equation}
Оскільки $S_v(y)$ не залежить від індексу~$i$, а $S_u(x)$ не залежить від індексу~$j$, можна розподілити ділення між обома сумами:
\begin{equation}\label{eq:m_L_distributed}
	L = \sum\limits_{i=0}^{n} \frac{u_i(x)}{S_u(x)} \cdot \sum\limits_{j=0}^{m} \frac{v_j(y)}{S_v(y)}\, f_{i,j}.
\end{equation}

\paragraph{Крок~4. Нормалізовані барицентричні коефіцієнти.}
Визначимо нормалізовані коефіцієнти:
\begin{equation}\label{eq:m_alpha_beta}
	\alpha_i(x) \stackrel{\text{def}}{=} \frac{u_i(x)}{S_u(x)} = \frac{\dfrac{w_i}{x - x_i}}{\sum\limits_{k=0}^{n} \dfrac{w_k}{x - x_k}}, \qquad
	\beta_j(y) \stackrel{\text{def}}{=} \frac{v_j(y)}{S_v(y)} = \frac{\dfrac{w_j}{y - y_j}}{\sum\limits_{k=0}^{m} \dfrac{w_k}{y - y_k}}.
\end{equation}
Безпосередньо з визначення випливає властивість нормалізації:
\begin{equation}\label{eq:m_norm_proof}
	\sum\limits_{i=0}^{n} \alpha_i(x) = \sum\limits_{i=0}^{n} \frac{u_i}{S_u} = \frac{1}{S_u} \sum\limits_{i=0}^{n} u_i = \frac{S_u}{S_u} = 1.
\end{equation}
Аналогічно $\sum_{j=0}^{m} \beta_j(y) = 1$. Іншими словами, $\alpha_i(x)$~--- це частка внеску $i$-го вузла у загальну суму ваг, і всі частки разом складають одиницю. Якщо $x$ збігається з вузлом $x_{i^*}$, покладаємо $\alpha_{i^*} = 1$, $\alpha_i = 0$ при $i \neq i^*$ (аналогічно для $\beta_j$).

У програмній реалізації нормалізація виконується одноразово на етапі
побудови вагових матриць: після обчислення всіх ваг
рядка знаходиться значення суми
$S_u(x_p)$, після чого кожен елемент рядка
ділиться на неї. Це повністю усуває операцію ділення з етапу
матричного множення.

\paragraph{Крок~5. Фінальна формула.}
Підставляючи~\eqref{eq:m_alpha_beta} у~\eqref{eq:m_L_distributed}, отримуємо:
\begin{equation}\label{eq:m_bilinear}
	L_{n,m}(x, y) = \sum\limits_{i=0}^{n}\sum\limits_{j=0}^{m} \alpha_i(x)\,\beta_j(y)\, f_{i,j}.
\end{equation}
Двовимірна барицентрична інтерполяція зведена до о зваженої суми значень $f_{i,j}$, де вага кожного пікселя є добутком двох одновимірних нормалізованих коефіцієнтів.

\paragraph{Крок~6. Матрична інтерпретація для однієї точки.}
Позначимо:
\begin{itemize}
	\item $\boldsymbol{\alpha}(x) = [\alpha_0(x),\; \alpha_1(x),\; \ldots,\; \alpha_n(x)]^{\!\top}$~--- вектор-стовпець розмірності $W$;
	\item $\boldsymbol{\beta}(y) = [\beta_0(y),\; \beta_1(y),\; \ldots,\; \beta_m(y)]^{\!\top}$~--- вектор-стовпець розмірності $H$;
	\item $\mathbf{F} \in \mathbb{R}^{H \times W}$~--- матриця кольорових значень вихідного зображення, $\mathbf{F}_{j,i} = f_{i,j}$.
\end{itemize}
Тоді формула~\eqref{eq:m_bilinear} записується як білінійна форма:
\begin{equation}\label{eq:m_bilinear_vec}
	L_{n,m}(x, y) = \boldsymbol{\beta}(y)^{\!\top}\,\mathbf{F}\, \boldsymbol{\alpha}(x),
\end{equation}
що обчислюється у два кроки:
\begin{enumerate}
	\item \textit{Інтерполяція вздовж $y$}:
	\begin{equation}\label{eq:m_pass1_vec}
		\mathbf{r}(y) = \mathbf{F}^{\!\top} \boldsymbol{\beta}(y) \in \mathbb{R}^{W}, \qquad
		r_i = \sum\limits_{j=0}^{m} \beta_j(y)\, f_{i,j}.
	\end{equation}
	\item \textit{Інтерполяція вздовж $x$}:
	\begin{equation}\label{eq:m_pass2_vec}
		L_{n,m}(x, y) = \boldsymbol{\alpha}(x)^{\!\top}\, \mathbf{r}(y) = \sum\limits_{i=0}^{n} \alpha_i(x)\, r_i.
	\end{equation}
\end{enumerate}

\paragraph{Крок~7. Узагальнення на всі цільові точки.}
Для цільового зображення розміром $W_2 \times H_2$ побудуємо дві \textit{вагові матриці}, рядки яких є нормалізованими барицентричними коефіцієнтами для кожного цільового пікселя:
\begin{equation}\label{eq:m_matrix_A}
	\mathbf{A} \in \mathbb{R}^{W_2 \times W}, \qquad \mathbf{A}_{p,i} = \alpha_i(x_p), \qquad p = 0, \ldots, W_2 - 1,
\end{equation}
\begin{equation}\label{eq:m_matrix_B}
	\mathbf{B} \in \mathbb{R}^{H_2 \times H}, \qquad \mathbf{B}_{q,j} = \beta_j(y_q), \qquad q = 0, \ldots, H_2 - 1,
\end{equation}
де $x_p = -\cos(\pi p / (W_2 - 1))$ та $y_q = -\cos(\pi q / (H_2 - 1))$. Кожен рядок обох матриць задовольняє умову нормалізації: сума елементів рядка дорівнює~$1$.

Елемент результуючої матриці $\mathbf{G}_{q,p} = L_{n,m}(x_p, y_q)$ відповідно до~\eqref{eq:m_bilinear} розраховується як:
\begin{equation*}
	\mathbf{G}_{q,p} = \sum\limits_{i=0}^{n}\sum\limits_{j=0}^{m} \mathbf{A}_{p,i}\, \mathbf{B}_{q,j}\, \mathbf{F}_{j,i}
	= \sum\limits_{i=0}^{n} \mathbf{A}_{p,i} \underbrace{\sum\limits_{j=0}^{m} \mathbf{B}_{q,j}\, \mathbf{F}_{j,i}} _{(\mathbf{B}\,\mathbf{F})_{q,i}}.
\end{equation*}
Внутрішня сума є елементом добутку матриць $\mathbf{B} \cdot \mathbf{F}$, а зовнішня~--- добутку з $\mathbf{A}^{\!\top}$. Отже, фінальне матричне рівняння має вигляд:
\begin{equation}\label{eq:m_result}
	\mathbf{G} = \mathbf{B}\, \mathbf{F}\, \mathbf{A}^{\!\top}.
\end{equation}

\paragraph{Поканальне представлення вхідних даних.}
Формула~\eqref{eq:m_result} визначена для матриці
$\mathbf{F}$. Оскільки колірне зображення формату RGB містить три
незалежні канали інтенсивності, у програмній реалізації матриця вхідних
даних декомпонується на три окремі матриці
$\mathbf{F}_R,\; \mathbf{F}_G,\; \mathbf{F}_B \in \mathbb{R}^{H \times W}$. Для кожного каналу матричне
множення виконується незалежно:
\begin{equation}\label{eq:m_result_rgb}
	\mathbf{G}_c = \mathbf{B}\, \mathbf{F}_c\, \mathbf{A}^{\!\top},
	\qquad c \in \{R, G, B\}.
\end{equation}
Такий підхід, на відміну від черезрядкового зберігання компонент, забезпечує послідовний доступ до пам'яті
всередині GEMM-ядра та сприяє автоматичній векторизації внутрішніх циклів
компілятором.

\paragraph{Програмна реалізація з бібліотекою Eigen.}
Для виконання матричного множення~\eqref{eq:m_result_rgb} використано
бібліотеку Eigen~--- header-only бібліотеку лінійної алгебри для C++, яка
реалізує апаратно оптимізований алгоритм GEMM
з автоматичним cache-blocking, використанням SIMD-інструкцій та
prefetch-директивами.

Вагові матриці $\mathbf{A}$ та $\mathbf{B}$ зберігаються у суцільних масивах типу \texttt{float} розмірністю $W_2 \times W$ та
$H_2 \times H$ відповідно. Для передачі їх до Eigen без копіювання даних
використовується механізм \texttt{Eigen::Map}, який створює матричний
об'єкт поверх існуючого масиву:
\begin{verbatim}
	Eigen::Map<MatRM> mWy(Wy.data(), H2, H);
	Eigen::Map<MatRM> mWx(Wx.data(), W2, W);
	Eigen::Map<MatRM> mFR(FR.data(), H, W);
\end{verbatim}
де \texttt{MatRM}~--- тип матриці у форматі
\texttt{Eigen::Matrix<float, Dynamic, Dynamic, RowMajor>}. Обчислення
результату для кожного каналу зводиться до єдиного виразу:
\begin{verbatim}
	MatRM mRR = (mWy * mFR) * mWx.transpose();
\end{verbatim}
Eigen автоматично створює проміжну матрицю
$\mathbf{T} = \mathbf{B} \cdot \mathbf{F}_c$ та виконує два послідовних
виклики оптимізованого GEMM-ядра.

Використання типу \texttt{float} замість \texttt{double} зменшує обсяг
даних удвічі, що подвоює ефективну пропускну здатність кеш-пам'яті та
дозволяє Eigen обробляти 8~елементів за один AVX-такт замість~4. Обчислення
самих вагових коефіцієнтів виконується у типі \texttt{double} для
збереження точності при діленні на малі різниці $(x - x_i)$, після чого
результат приводиться до \texttt{float}.

\paragraph{Побудова вагових матриць.}
Побудова матриць $\mathbf{A}$ та $\mathbf{B}$ виконується з використанням
усіх доступних апаратних потоків: діапазон цільових індексів рівномірно
розподіляється між потоками за допомогою пулу \texttt{std::thread}.
Спочатку всі доступні потоки будують матрицю $\mathbf{B}$, після завершення~- всі потоки будують
матрицю~$\mathbf{A}$.

На етапі побудови застосовано ті самі оптимізації, що й у злитому
алгоритмі: детекція точного збігу з вузлом за $O(1)$ на основі
цілочисельної арифметики, чергування знаків
та попередня нормалізація ваг, що усуває операцію ділення з етапу
GEMM-множення.

\paragraph{Конвертація вхідних даних.}
Вхідне зображення \texttt{QImage} конвертується у формат
\texttt{Format\_RGB32}, після чого колірні компоненти кожного пікселя
витягуються безпосередньо з \texttt{QRgb}-значень та записуються у три
окремі суцільні \texttt{float}-масиви $\mathbf{F}_R$, $\mathbf{F}_G$,
$\mathbf{F}_B$. Такий підхід усуває проміжні конвертації формату, що економить декілька повних проходів по масиву пікселів.

Результуючі значення після GEMM-множення обмежуються стандартним
колірним діапазоном $[0, 255]$ за допомогою \texttt{std::clamp} та
записуються безпосередньо у вихідний \texttt{QImage} формату
\texttt{Format\_RGB32} без проміжних матриць у цілочисельному форматі.

\paragraph{Обчислювальна складність.}
Побудова вагових матриць потребує $\mathcal{O}(W_2 \cdot W)$ операцій
для $\mathbf{A}$ та $\mathcal{O}(H_2 \cdot H)$ для~$\mathbf{B}$.
Матричне множення $\mathbf{B} \cdot \mathbf{F}_c$ виконується за
$\mathcal{O}(H_2 \cdot W \cdot H)$ операцій, а множення
$\mathbf{T}_c \cdot \mathbf{A}^{\!\top}$~--- за
$\mathcal{O}(H_2 \cdot W_2 \cdot W)$. Для трьох каналів загальна
асимптотична складність алгоритму становить:
\begin{equation}
	\mathcal{O}\bigl(3 \cdot W \cdot H_2 \cdot (H + W_2)\bigr)
	= \mathcal{O}\bigl(W \cdot H_2 \cdot (H + W_2)\bigr).
\end{equation}

Асимптотична складність матричної форми збігається зі складністю злитого
алгоритму. Однак матрична реалізація з використанням Eigen досягає
суттєво вищої практичної швидкодії завдяки оптимізованому GEMM-ядру:
cache-blocking забезпечує роботу з даними у кеш-пам'яті, SIMD-інструкції
обробляють 8~значень \texttt{float} за такт, а prefetch-директиви
зменшують затримку доступу до RAM. У сукупності ці фактори
забезпечують пропускну здатність GEMM, що наближається до теоретичної
пікової продуктивності процесора.